% ============================================================
% SECTION 1 — 'Ali's Struggles (Chapter 1)
% ============================================================
\section{'Ali's Struggles}

\begin{frame}{Justice Has a Cost}
\bigquote{Are you asking me to gain victory and success in politics at the price of oppression and destroying the rights of powerless people?!}{'Ali, refusing special subsidies to influential allies}
\begin{itemize}[itemsep=5pt]
  \item 'Ali initially refused the caliphate --- he foresaw the trouble ahead
  \item 'Uthman's assassination created a new kind of enemy: intelligent, non-religious, dressed in Islamic language
  \item Mu'awiyah displayed 'Uthman's blood-stained shirt in Damascus --- manufactured outrage, no genuine kinship claim
  \item 'Ali's own camp pushed him to be ``more tactful'' like Mu'awiyah --- he refused: deception is available to him too, but ``immoral, evil, wicked and blasphemous''
\end{itemize}
\end{frame}

\begin{frame}{The Structural Cause of the Kharijite Crisis}
\begin{itemize}[itemsep=5pt]
  \item Prophet: 13 pre-Hijrah years training a small core --- jihad \emph{forbidden}, even under torture
  \item Later expansion under 'Uthman: rapid conquest, no equivalent formation period
  \item Result: converts with Islam's outward ``crust'' but not its substance
\end{itemize}
\bigquote{They do not know what the Qur'an is. They do not understand the meaning of the Qur'an.}{'Ali, on the future Kharijites}
\delivernote{This is the chapter's diagnostic claim: the Kharijites are a predictable product of skipping the Prophet's slow formation method, not simply bad people appearing from nowhere.}
\end{frame}

\begin{frame}{Qur'ans on Spears}
\begin{itemize}[itemsep=5pt]
  \item Siffin: 'Ali's side on the verge of winning
  \item Mu'awiyah's forces raise Qur'ans on spear-points ('Amr ibn al-'As's idea)
  \item 'Ali's own pietist faction refuses to keep fighting
\end{itemize}
\bigquote{They have not brought forward the Qur'an itself but its papers and cover... I am your imam. I am your speaking Qur'an.}{'Ali, objecting}
\examplenote{The trap}{20{,}000 threatened to kill 'Ali himself unless the near-victorious commander Malik al-Ashtar was recalled from the field. He was.}
\end{frame}

\begin{frame}{Arbitration $\rightarrow$ Kharijite Doctrine}
\begin{itemize}[itemsep=5pt]
  \item Faction forces Abu Musa al-Ash'ari onto 'Ali as arbitrator, overriding his choice of Malik or Ibn 'Abbas
  \item Abu Musa is outmaneuvered by 'Amr ibn al-'As
  \item Rather than admit the mistake, the faction declares \emph{arbitration itself} blasphemy (Qur'an 6:57) and demands 'Ali repent
\end{itemize}
\bigquote{Arbitration is no blasphemy... I shall never confess to a sin I have never committed.}{'Ali}
\delivernote{Ask: why is ``we were wrong to force this ceasefire'' a much harder conclusion to reach than ``arbitration itself is heresy''? What does that reveal about how doctrinal extremism forms?}
\end{frame}

\begin{frame}{Kharijite Doctrine in Practice}
\begin{itemize}[itemsep=5pt]
  \item 'Uthman, 'Ali, and Mu'awiyah all declared non-believers for accepting arbitration
  \item A single major sin $\rightarrow$ expulsion from Islam entirely
  \item In practice: killed a companion of the Prophet and his pregnant wife for refusing to disown 'Ali
  \item Rebuked a man for eating one unattended date as ``intruding on your Muslim brother's wealth''
\end{itemize}
\end{frame}

\begin{frame}{Nahrawan --- and What Came After}
\begin{itemize}[itemsep=5pt]
  \item 12{,}000 Kharijites confronted; 8{,}000 repented under 'Ali's raised banner; 4{,}000 died in battle
  \item ``It was I and I alone who removed the eye of this revolt'' --- 'Ali claims a unique moral authority: an ordinary believer's hand would shiver at fighting people who \emph{look} more ascetic than themselves
  \item Ibn Muljam survives Nahrawan --- his father and brother killed there is later exploited by Qatam, who demands 'Ali's assassination as her final dowry condition
\end{itemize}
\bigquote{After me, do not kill the Kharijites anymore, because whoever seeks truth and commits a mistake is not the same as the one who seeks falsehood from the beginning.}{'Ali, his final instruction on the sect that killed him}
\delivernote{Modern parallel worth telling here: the Iqbal/Ja'far-Sadiq rumor. A couplet condemning two colonial-era traitors coincidentally named Ja'far and Sadiq was misread as cursing the Imams. Mutahhari's response: ``Investigate! Investigate... do not say anything before you have done your research.''}
\end{frame}

\qaframe{'Ali's Struggles}{%
  \qa{State the distinction between ignorance and hypocrisy that this chapter builds, and argue whether it holds up under pressure.}
     {Ignorance (Kharijites) is sincere but dangerous and needs patience, correction, force only as last resort. Hypocrisy (Mu'awiyah's circle) is calculated and needs exposure, not correction. Nahrawan tests this: 8,000 of 12,000 do repent when confronted -- real evidence the distinction matters -- but the remaining 4,000 still needed exactly the force a hypocrite would require.}
  \qa{Trace the arc from the training gap, to Kharijite doctrine, to 'Ali's final instruction. What does it argue about how extremism is created and answered?}
     {Skipping the Prophet's slow formation method produces hollow piety; battlefield pressure hardens that hollowness into a doctrine licensing violence; 'Ali's answer, even dying, is not indiscriminate retaliation but a permanent instruction to keep distinguishing sincere error from deliberate falsehood. Prevention, not punishment, is the real fix.}
  \qa{How does the Iqbal rumor anticipate a concern that resurfaces in Chapter 8, Part 2?}
     {Chapter 8 repeatedly warns against mistaking a claimant, a name-coincidence, or a rumor for the real thing when it comes to Mahdi claims (Muhammad ibn 'Ijlan's mistaken allegiance is the direct parallel). The Iqbal story plants that same discipline decades earlier.}
}
